% =========================================================
% Square Bracketing and Nonlinear Approximation
% =========================================================

\section{Square Bracketing and Nonlinear Approximation}
\label{sec:square-bracketing-nonlinear-approximation}

Linear approximation controls distances on the number line.  Nonlinear approximation asks for rational inputs whose images under nonlinear operations, such as squaring, bracket or approximate a target value.

\begin{tcolorbox}[colback=thmbox, colframe=thmborder, arc=2pt,
  left=6pt, right=6pt, top=4pt, bottom=4pt,
  title={\small\textbf{Toolkit --- Square Approximation}},
  fonttitle=\small\bfseries]
The basic square-approximation moves are:
\[
a^2<r<b^2,
\]
then refine the bracket or choose a fine rational grid and take the first grid point whose square crosses $r$.
\end{tcolorbox}

\begin{tcolorbox}[colback=thmbox, colframe=thmborder, arc=2pt,
  left=6pt, right=6pt, top=4pt, bottom=4pt,
  title={\small\textbf{Theorem (Rational Square Bracketing)}},
  fonttitle=\small\bfseries]
\begin{theorem}[Rational Square Bracketing]
\label{thm:rational-square-bracketing}
For every $r\in\mathbb{Q}$ with $r>0$, there exist $a,b\in\mathbb{Q}$ with $a>0$ and $b>0$ such that
\[
a^2<r<b^2.
\]
\end{theorem}
\end{tcolorbox}

\begin{remark*}[Proof idea]
Use
\[
a=\frac{r}{r+1},\qquad b=r+1.
\]
\end{remark*}

\begin{remark*}[Status]
Discussed and proof-checked in conversation; final proof file is not present in the uploaded notes zip.
\end{remark*}

\begin{tcolorbox}[colback=thmbox, colframe=thmborder, arc=2pt,
  left=6pt, right=6pt, top=4pt, bottom=4pt,
  title={\small\textbf{Lemma (First Grid Point Crossing)}},
  fonttitle=\small\bfseries]
\begin{lemma}[First Grid Point Crossing]
\label{lem:first-grid-point-crossing}
Let $r\in\mathbb{Q}$ with $r>0$ and let $q\in\mathbb{N}$ with $q\ge 1$. If $K\in\mathbb{N}$ satisfies $K^2>r$, then the set
\[
A=\{m\in\mathbb{N}:m^2>rq^2\}
\]
is nonempty and therefore has a least element.
\end{lemma}
\end{tcolorbox}

\begin{remark*}[Status]
Missing final proof. This uses the well-ordering principle and the witness $Kq\in A$.
\end{remark*}

\begin{tcolorbox}[colback=thmbox, colframe=thmborder, arc=2pt,
  left=6pt, right=6pt, top=4pt, bottom=4pt,
  title={\small\textbf{Theorem (Rational Square Approximation from Above)}},
  fonttitle=\small\bfseries]
\begin{theorem}[Rational Square Approximation from Above]
\label{thm:rational-square-approximation-from-above}
Let $r\in\mathbb{Q}$ with $r>0$, and let $n\in\mathbb{N}$ with $n\ge 1$. Then there exists $s\in\mathbb{Q}$ with $s>0$ such that
\[
r<s^2<r+\frac1n.
\]
\end{theorem}
\end{tcolorbox}

\begin{remark*}[Proof idea]
Choose a large grid denominator $q$, then let $p$ be the least natural number such that $p^2>rq^2$.  Set $s=p/q$ and bound the overshoot using
\[
p^2=(p-1)^2+2p-1.
\]
\end{remark*}

\begin{remark*}[Status]
Discussed in conversation; final proof file is not present in the uploaded notes zip.
\end{remark*}

\begin{tcolorbox}[colback=thmbox, colframe=thmborder, arc=2pt,
  left=6pt, right=6pt, top=4pt, bottom=4pt,
  title={\small\textbf{Theorem (Rational Square Approximation from Below)}},
  fonttitle=\small\bfseries]
\begin{theorem}[Rational Square Approximation from Below]
\label{thm:rational-square-approximation-from-below}
Let $r\in\mathbb{Q}$ with $r>0$, and let $n\in\mathbb{N}$ with $n\ge 1$. If $r-1/n>0$, then there exists $s\in\mathbb{Q}$ with $s>0$ such that
\[
r-\frac1n<s^2<r.
\]
\end{theorem}
\end{tcolorbox}

\begin{remark*}[Status]
Missing proof.
\end{remark*}
